\begin{table}[!t]
    \caption{
        大规模仿真中使用的模型组合。
    }
    \label{table:xl-model-setup}

    \footnotesize
    \begin{tabular}{llcccc}
    \toprule
    \textbf{名称} & \textbf{模型配置} & \textbf{DP} & \textbf{TP} & \textbf{PP} & \textbf{GPU 数} \\
    \midrule
    VLM-XL-8k & ViT 22B + GPT 175B & 128 & 8 & 8 & 8192 \\
    VLM-XL-16k & ViT 22B + GPT 175B & 128 & 8 & 16 & 16384 \\
    \midrule
    T2V-XL-3k & Qwen2 72B + DiT 30B & 96 & 8 & 4 & 3072 \\
    T2V-XL-6k & Qwen2 72B + DiT 30B & 96 & 8 & 8 & 6144 \\
    \bottomrule
    \end{tabular}
\end{table}

\section{结论}

由于流水线阶段不均衡和训练数据动态性，高效训练大多模态模型颇具挑战。
本文提出 {\sys}，通过自适应模态感知划分和高效的流水线调度搜索，
解决 LMM 训练中的动态不均衡。
实验结果表明，相比现有最先进训练系统，{\sys} 可将训练吞吐量最高提升 97.3\%。
